%% footnotenosplit.tex -- documentation for the footnotenosplit package
%%
%% Copyright (c) 2026 Munehiro Yamamoto <munepixyz@gmail.com>
%% License: LPPL 1.3c or later
%%
%% Build: pdflatex footnotenosplit.tex (twice)
\documentclass{article}
\usepackage[margin=28mm]{geometry}
\usepackage{xcolor}
\usepackage[colorlinks, allcolors=blue!50!black]{hyperref}
\newcommand*\fnsversion{v1.1 (2026/11/17)}
\newcommand*\cs[1]{\texttt{\char`\\#1}}
\newcommand*\pkg[1]{\textsf{#1}}

\title{The \pkg{footnotenosplit} package\\
  \large Suppress orphan footnote rules when footnotes are not allowed to split}
\author{Munehiro Yamamoto\\
  \url{https://github.com/munepi/footnotenosplit}}
\date{\fnsversion}

\begin{document}
\maketitle

\section{The problem}

Setting \cs{interfootnotelinepenalty}\texttt{=10000} keeps footnote
text from being split across pages: a footnote that does not fit on the
current page is moved to the next page as a whole.  This policy is
common in book production, but it has two side effects.

\begin{enumerate}
\item The original page may end with a \emph{lone} \cs{footnoterule}
  (plus the \cs{skip}\cs{footins} space) with \emph{no footnote text
  below it}.
\item The footnote moved to the next page gets only \cs{splittopskip}
  ($=$ \cs{footnotesep}) above its first line, so it sits closer to the
  rule than a normal footnote does.
\end{enumerate}

\noindent
Schematically, the bottom of the page where the footnote call sits:

\begin{center}
\newcommand*\pagebottomdemo[2]{%
  \begin{minipage}[t]{0.46\linewidth}
    \centering
    \fbox{\begin{minipage}[t]{0.88\linewidth}
      \footnotesize
      \ldots{} filler text filler text filler text\\
      the reference line sits here\textsuperscript{2}\\[4pt]
      #1%
    \end{minipage}}\\[4pt]
    {\footnotesize(#2)}%
  \end{minipage}}%
\pagebottomdemo{\rule{0.6\linewidth}{0.4pt}\\[8pt]\mbox{}}%
  {without \pkg{footnotenosplit}: a lone rule}%
\hfill
\pagebottomdemo{\mbox{}\\[8pt]\mbox{}}%
  {with \pkg{footnotenosplit}: no rule}%
\end{center}

\paragraph{Why it happens.}
When the footnote insertion begins with vertical material---typically a
leading \cs{kern} at the top of \cs{@makefntext}, as in many
book-production classes---and the page bottom has stretch
(\cs{raggedbottom}), \TeX's page builder ``splits'' the unsplittable
insertion at its very top.  The page then keeps an \emph{empty}
\cs{box}\cs{footins} that contains only that leading kern, and the
\LaTeX{} output routine, which only tests \cs{ifvoid}\cs{footins},
duly typesets \cs{footnoterule} above the empty box.  The phenomenon is
a property of \TeX's page builder and the \LaTeX{} output routine: it
occurs with any engine (pdf\LaTeX, Lua\LaTeX, Xe\LaTeX, (u)p\LaTeX)
and any language.

\section{What the package does}

\begin{enumerate}
\item At the start of the output routine, if
  $\cs{ht}\cs{footins} < \cs{footnotesep}$ (i.e.\ the footnote box
  holds no actual footnote line), \cs{box}\cs{footins} is discarded,
  so neither \cs{footnoterule} nor \cs{skip}\cs{footins} is emitted.
  The footnote itself is typeset normally on the next page.
\item \cs{splittopskip} inside \cs{@footnotetext}
  (\cs{H@@footnotetext} when \pkg{hyperref} is loaded) is patched to
  $\cs{footnotesep} + \cs{footnotetopkern}$, so that the moved footnote
  gets the same visual space below the rule as a normal footnote.
\end{enumerate}

\paragraph{Tagged PDF.}
The package also works with \verb|\DocumentMetadata{tagging=on}|
(pdf\LaTeX{} and Lua\LaTeX).  There, the \LaTeX{} kernel's tagging
code replaces \cs{@footnotetext}, so the package sets the adjusted
\cs{splittopskip} in the \texttt{fntext/para} hook instead of patching
\cs{@footnotetext}.  The box discarded in item~1 is empty in tagged
documents as well, so the structure tree is not affected.

\section{Usage}

\begin{verbatim}
\usepackage{footnotenosplit}

% If your class's \@makefntext starts with "\kern <dimen>",
% set \footnotetopkern to that dimension and replace the \kern:
\footnotetopkern=2pt
\makeatletter
\renewcommand\@makefntext[1]{\kern\footnotetopkern ... #1}
\makeatother
\end{verbatim}

\noindent
The dimen register \cs{footnotetopkern} (default: \texttt{0pt}) tells
the package how much vertical material \cs{@makefntext} puts above the
first footnote line; it is added to \cs{splittopskip} for footnotes
moved to the next page.

\paragraph{Option \texttt{nopenalty}.}
By default the package sets
\cs{interfootnotelinepenalty}\texttt{=}\cs{@M}.  If your class already
sets it, or you want to manage it yourself, load the package with
\verb|\usepackage[nopenalty]{footnotenosplit}|.

\section{Implementation notes}

The empty-box test is engine-independent, but where to hook it is not,
so \pkg{footnotenosplit.sty} dispatches to one of two implementation
files (do not load them directly):

\begin{description}
\item[\pkg{footnotenosplit-latex.sty}] for \LaTeX{} with the
  socket-based \cs{@makecol} (\LaTeX2e 2025-06-01 or later; pdf\LaTeX,
  Lua\LaTeX, Xe\LaTeX, \dots).  It adds the test to the
  \texttt{build/column/before} hook, executed at the very beginning of
  \cs{@makecol}; the subsequent \cs{@outputbox@appendfootnotes} then
  sees a void \cs{footins} and emits neither \cs{footnoterule} nor
  \cs{skip}\cs{footins}.
\item[\pkg{footnotenosplit-ptex.sty}] for (u)p\LaTeX, whose
  \texttt{plcore.ltx} overrides \cs{@makecol} with a pre-socket
  definition.  It patches that \cs{@makecol} directly (via
  \pkg{etoolbox}'s \cs{patchcmd}).
\end{description}

On an older \LaTeX{} kernel (before 2025-06-01) without the socket
mechanism, the package issues a warning and only applies the
\cs{splittopskip} patch.

\section{Tests}

The repository ships \texttt{test-latex.tex} (pdf\LaTeX/Lua\LaTeX/%
Xe\LaTeX) and \texttt{test-ptex.tex} ((u)p\LaTeX{} $+$
\texttt{dvipdfmx}), which reproduce the lone rule and verify its
suppression, plus \texttt{test-latex-show.tex} and
\texttt{test-ptex-show.tex}, which dump the patched macros to the log.
The pdf\LaTeX/Lua\LaTeX{} tests also have tagged PDF variants
(\verb|\def\TAGGING{1}|, optionally with \verb|\def\HYPERREF{1}|).
See the comments in each file.

\section{License}

This package is licensed under the conditions of the \LaTeX{} Project
Public License, either version~1.3c of this license or any later
version.

\medskip
\noindent
Copyright \copyright{} 2026 Munehiro Yamamoto

\end{document}
